\documentclass{article}

\textwidth=6.5in
\textheight=8.5in
\voffset=-54pt
\oddsidemargin=0in
\evensidemargin=0in
\setlength{\parskip}{8pt}
\setlength{\parindent}{0pt}

\usepackage{amsmath, amssymb}
\usepackage{ifthen}

\usepackage[inline]{enumitem}
\setlist{topsep=0pt, parsep=8pt}

\usepackage[rgb,table]{xcolor}\convertcolorsUfalse
\usepackage{graphicx}

% change hyperref options
\PassOptionsToPackage{hyphens}{url}
\usepackage[bookmarks,colorlinks,linkcolor=black,citecolor=blue,pdfstartview=FitH,urlcolor=blue]{hyperref}
\hypersetup{pdfpagemode=UseNone}
\usepackage{bookmark}

\usepackage{graphicx}
\usepackage{multicol}
\usepackage{framed}
\usepackage{array}

\usepackage{icomma}

\usepackage{soul}
\usepackage[normalem]{ulem}

\usepackage{pgfplots}
\pgfplotsset{compat=1.18}  
\pgfplotscreateplotcyclelist{mycolor}{black, blue, red, green!70!black, magenta, purple, cyan, orange }  
\pgfplotsset{
  disabledatascaling,
  cycle list=tikzcolor, 
  axis lines=middle,
  every axis plot/.style={no marks, thick},
  every axis/.style={
    enlarge x limits=false,
    enlarge y limits=auto,
    xlabel style={at={(current axis.right of origin)},anchor=west},
    ylabel style={at={(current axis.above origin)},anchor=south},
    % extra x and y ticks for asymptotes
    extra x tick style={grid=major, grid style={dashed,black}},
    extra y tick style={grid=major, grid style={dashed,black}},
    /pgf/number format/1000 sep={},
  },    
  tick label style = {font=\footnotesize, fill=\currentbgcolor, inner sep=0pt},    
  xticklabel shift=1pt,    
  scaled ticks=false,
  scaled y ticks=false,
  unbounded coords=jump,
  samples=101,
  minor tick length=0,
  scatter/classes={
    open={mark=*,draw=black,fill=white, mark size=2, thin},
    closed={mark=*,draw=black,fill=black, mark size=2, thin}
  },
  width=3in,
}
\pgfplotsset{open/.style={only marks, mark=*, fill=white, thin, mark size=2}}
\pgfplotsset{closed/.style={only marks, mark=*, draw=black, fill=black,
    thin, mark size=2}}
\pgfplotsset{labeled/.style={nodes near coords, only marks, mark=*, draw=black, fill=black, thin, mark size=2, point meta=explicit symbolic}}

\pgfplotsset{gridgraph/.style={clip=false, axis line style={shorten
      >=-8pt}, xlabel style={xshift=8pt}, ylabel style={yshift=8pt},
    enlargelimits=false, grid=major, tick style={black}}} 

\usepackage{tikz}
\usetikzlibrary{
  matrix,%
  % enables the snake paths
  decorations.pathmorphing,%
  % enables bigger arrows
  decorations.markings,%
  decorations.pathreplacing,%
  decorations.text,%
  calc,%
  patterns,%
  shapes.geometric,%
  arrows,
  decorations.shapes,
  positioning,
  plotmarks,
  intersections
}
\tikzset{>=latex} % sets default arrow type in tikz
\tikzset{font=\normalsize}
\tikzset{mark size=1.5pt, mark options=thin}
\tikzset{pin distance=4pt,
  every pin edge/.style={<-, thin, shorten <= -2pt}}

\interfootnotelinepenalty=10000
\renewcommand{\thefootnote}{(\arabic{footnote})}

\usepackage{multirow, bigdelim}

% change to \usepackage[sols]{handout} to see solutions
\usepackage[sols, grading, notes, workshop]{handout}

\newcommand\iscurrentenumi[1]{%
  \ifthenelse{\equal{#1}{\theenumi}}{}{#1}}
\setenumerate[1]{ref=\#\arabic*}
\setenumerate[2]{ref=\protect\iscurrentenumi{\theenumi}(\alph*)}
\newcommand\enumiiprefix[2]{%
  \ifthenelse{{\equal{#1}{\theenumi}}\AND{\equal{#2}{(\alph{enumii})}}}{}{}%
  \ifthenelse{{\equal{#1}{\theenumi}}\AND{\NOT{\equal{#2}{(\alph{enumii})}}}}{#2}{}%
  \ifthenelse{\equal{#1}{\theenumi}}{}{#1#2}%
}
\setenumerate[3]{ref=\protect\enumiiprefix{\theenumi}{(\alph{enumii})}\roman*}

\makeatletter

% underline links               
\let\@href=\href
\renewcommand{\href}[2]{\@href{#1}{\ul{#2}}}

\newdimen\SOUL@dimen %new
\def\SOUL@ulunderline#1{{%
    \setbox\z@\hbox{#1}%
    % \dimen@=\wd\z@
    \SOUL@dimen=\wd\z@ %new
    \dimen@i=\SOUL@uloverlap
    \advance\SOUL@dimen2\dimen@i %\dimen@ exchanged too
    \rlap{%
      \null
      \kern-\dimen@i
      % \SOUL@ulcolor{\SOUL@ulleaders\hskip\dimen@}%
      \SOUL@ulcolor{\SOUL@ulleaders\hskip\SOUL@dimen}% new
    }%
    \unhcopy\z@
  }}

\newcommand\calculator{
  \begin{tikzpicture}[x=0.15ex, y=0.15ex, baseline=1.5pt]
    \begin{scope}[thick]
      \draw (0, 0) -- (10, 0) -- (10, 14) -- (0, 14) -- cycle;
      \foreach \x in {10/3, 20/3}
      {
        \draw (\x, 0) -- (\x, 10);
      }
      \foreach \y in {10/3, 20/3, 10}
      {
        \draw (0, \y) -- (10, \y);
      }
    \end{scope}
  \end{tikzpicture}
}
\newcommand{\calcitem}{\refstepcounter{enum\romannumeral\@enumdepth}%
\item[\calculator \csname labelenum\romannumeral\@enumdepth\endcsname]}

\makeatother

\definecolor{purple}{rgb}{0.5,0,1.0}

\newcommand{\doedfinity}[2][\newpage]{
  \begin{problem}[#1]
    Do the Edfinity assignment ``Problem Set #2''.

    We encourage you to write down any scratch work here, as well as
    things you want your future self to remember. Nothing you write
    here will be graded; it's just for you!
  \end{problem}
}

\newcommand{\psrnewpage}{\newpage}
\newcommand{\psreflection}[2][]{

  \ifthenelse{\not\equal{#1}{nocite}}{
    \begin{enumerate}[resume]
    \item
      \begin{problem}[1in]
        Please cite any help you got on this problem set:
        \begin{itemize}[nosep]
        \item If you worked with other students, please list their names here.
        \item If you used resources other than those on our Canvas site, please list them here.
        \end{itemize}
      \end{problem}

    \end{enumerate}
  }

  \probonly{%
    #2

    \psrnewpage

    \emph{Reflection space.} It's a good habit to take a few minutes at the end of each pset to reflect on what you learned and jot down notes to your future self. Here are a few reflection questions to get you started. Nothing you write here will be graded; it's just for you!
    \begin{itemize}
    \item Take a few minutes to look back at the problems. How could you explain to someone else how to do each problem? What was your strategy, and how did you know to use that strategy? If there are problems where you don't feel able to answer this, write down which ones they are. Come back to them in a few days when we post the homework solutions!

      \vspace{1.5in}

    \item Take a look back at the learning objectives on today's worksheet; how are you feeling about each one? If there are ones you know you need more practice with, write those down here.

      \vspace{1.5in}

    \item What did you find most challenging about this material? Do you have any lingering questions about it? If so, write them down here so you don't forget them. At some point, stop by office hours or MQC to ask!

      \vfill
    \end{itemize}      
  }
}

\usepackage{fancyhdr}
\lhead{\textcolor{white}{This content is protected and may not be shared, uploaded, or distributed.}}
\rhead{}
\rfoot{\vspace*{\fill}\footnotesize\textcopyright{} \emph{Harvard Math Ma}}
\renewcommand{\headrulewidth}{0pt}
\pagestyle{fancy}
\begin{document}

\begin{center}
  \large\textsc{Problem Set 24 - Applications of Exponentials}\vspace{.5in}\\
  \begin{problemonly}
   \large Name: \rule{5cm}{0.15mm}\\
   \end{problemonly}
\end{center}

\begin{enumerate}
\item  \note{
    \begin{itemize}[nosep]
    \item Pick the picture that shows the correct relationship between the graphs $y = 2^x$ and $y = 3^x$. 
    \item Model compound interest
    \item Model something with a half-life and something decreasing by a certain percent per decade, and express both as $C \cdot b^t$.
    \item Given a formula for an exponential (like $P_0 (1.007)^t$), identify the percent increase / decrease per year.
    \item Match exponentials and vertical shifts to their graphs.      
    \end{itemize}
  }  

  \doedfinity{24}

\item % 9.3.3 modified

  A population of bacteria is growing exponentially. At 7:00 am, the mass of the population is 12 mg. Five hours later, it is 14 mg.

  \begin{enumerate}

  \item

    \begin{grading}
      0.5 points
    \end{grading}

    \begin{problem}[0.25in]
      Fill in the blank: every 5 hours, the bacteria's mass is multiplied by $\fitb{0.25in}{\frac{7}{6}}$.
    \end{problem}

  \item

    \begin{grading}
      2 points:
      \begin{itemize}[nosep]
      \item 0.5 for having the population at time 0 be 12
      \item 0.5 for having $\left( \frac{7}{6} \right)^\text{something}$. (If they didn't get $\frac{7}{6}$ in (a), they should instead use whatever value they got in (a).)
      \item 1 for having $t / 5$ in the exponent
      \end{itemize}

    \end{grading}

    \begin{problem}[\vfill]
      Write a formula for the bacteria's mass $t$ hours after 7:00 am. 
    \end{problem}

    \begin{solution}
      $M = 12\left(\frac{7}{6}\right)^{t/5}$.
    \end{solution}

  \calcitem

    \begin{grading}
      1.5 points: 1 for the exact expression $\left( \dfrac{7}{6} \right)^{24/5} - 1$ and 0.5 for the approximation expressed as a percent.
    \end{grading}

    \begin{problem}[\stretch{0.5}]
      By what percent is the bacteria's mass increasing each day? Write an exact expression, and use a calculator to give the approximate value as a percent. 
    \end{problem}

    \begin{solution}
      Each 24 hours, the mass is multiplied by $\left(\frac{7}{6}\right)^{24/5}$, so the
      percent change is $\left(\frac{7}{6}\right)^{24/5}-1$, which is approximately
      $1.1$ or $110\%$.
    \end{solution}

  \item

    \begin{grading}
      1 point
    \end{grading}

    \begin{problem}[\nstretch{0.5}]
      The \uline{doubling time} is the amount of time it takes for an exponentially increasing quantity to double in size. Is the doubling time for this bacteria population more or less than 24 hours? Explain briefly how you know. 
    \end{problem}

    \begin{solution}
      The doubling time for this bacteria is less than 24 hours. We know this because
      each 24 hours, the mass of bacteria increases by approximately $106\%$, more than
      doubling.
    \end{solution}

  \end{enumerate}

\calcitem \label{exponential-misconceptions}

  \emph{This problem highlights some common misconceptions about exponential functions.}

  \begin{enumerate}

  \item The price of eggs is increasing by 10\% every year. Luke says, ``That means that, in 5 years, the price of eggs will have increased by 50\%.'' But Luke is wrong!

    \begin{enumerate}
    \item

      \begin{grading}
        2 points
      \end{grading}

      \begin{problem}[\stretch{1.5}]
        By what percent will the price of eggs actually have increased after 5 years? Show your work.
      \end{problem}

      \begin{solution}
        If the current price of eggs is $P_0$, the price after $t$ years is $P_0 (1.1)^t$. So, in 5 years, the price will be $P_0 (1.1)^5 = 1.61051P_0$, which means the price will have increased by about 61\%.
      \end{solution}

    \item

      \begin{grading}
        2 points
      \end{grading}

      \begin{problem}[\vfill]
        Give Luke an intuitive explanation of why he's wrong. Is the percent increase over 5 years more or less than 50\%, and why does that make sense? 
      \end{problem}

      \begin{solution}
        Each year, the price increases by 10\% from the previous year's price, not from the initial price $P_0$. Since each year's price is higher than the last, each year's 10\% increase is also a larger amount than the last. Thus, the increase over 5 years  is more than a $50\%$ increase from $P_0$. 
      \end{solution}

    \end{enumerate}

  \item According to \href{https://canvas.harvard.edu/files/21163133/download?download_frd=1}{this article}, eating vegetables can improve mental health. In particular, % https://www.cambridge.org/core/services/aop-cambridge-core/content/view/06F5410553CF2C3849AAB0D9CE56E9B5/S0007114518000697a.pdf/div-class-title-fruit-and-vegetable-consumption-and-risk-of-depression-accumulative-evidence-from-an-updated-systematic-review-and-meta-analysis-of-epidemiological-studies-div.pdf
    ``a meta-analysis of 18 studies found that for every 100 grams of vegetables consumed [daily], depression risk dropped by 3 percent.'' Will says, ``That means that, if you consume an additional 500 grams of vegetables daily, your depression risk will drop by 15\%.'' But Will is incorrect.

    \begin{enumerate}
    \item

      \begin{grading}
        2 points
      \end{grading}

      \begin{problem}[\nstretch{1.5}]
        By what percent would your depression risk actually decrease? Please show your work. 
      \end{problem}

      \begin{solution}
        If your baseline depression risk is $D_0$, then according to this, your depression
        risk if you consume $x$ hundred additional grams of vegetables daily 
        is $D_0(0.97)^x$. So
        if you consume $500$ additional grams of vegetables daily, your depression risk
        is $D_0(0.97)^5$, or approximately $0.859D_0$. This is a percent change
        of $(0.97)^5-1 \approx =0.141$, or approximately $-14.1\%$.
      \end{solution}

    \item

      \begin{grading}
        2 points
      \end{grading}

      \begin{problem}[\vfill]
        Give Will an intuitive explanation of why he's wrong. Is the percent decrease from eating an additional 500 grams of vegetables daily more or less than 15\%, and why does that make sense? 
      \end{problem}

      \begin{solution}
        For every additional 100 grams of vegetables you consume daily, your depression
        risk decreases by $3\%$ compared to your current depression risk, not your
        baseline depression risk. Since your depression risk goes down as you consume
        more vegetables, every $3\%$ decrease is a smaller amount compared to before.
        Thus, the percent decrease from consuming an additional
        500 grams of vegetables is less than $15\%$.
      \end{solution}

    \end{enumerate}
  \end{enumerate}

\item 

  Medieval legends claim that Camelot was really the city of Winchester in England. The \href{http://commons.wikimedia.org/wiki/File:Winchester_-_Table_ronde_du_roi_Arthur.JPG}{Winchester Round Table}, an enormous table hanging in Winchester Castle, has often been identified as King Arthur's Round Table. According to legend, King Arthur lived in the 5th century.

  In 1976, scientists decided to use a method called radiocarbon dating to determine the age of the Winchester Round Table. This method relies on the fact that plants absorb two types of carbon; all living plants contain approximately the same ratio of radioactive carbon-14 to non-radioactive carbon-12. When a plant dies, its carbon-14 decays, but its carbon-12 doesn't. Therefore, after the plant dies, its ratio of carbon-14 to carbon-12 gradually declines; this ratio has a half-life of 5730 years.

  \begin{enumerate}
  \item

    \begin{grading}
      2 points
    \end{grading}

    \begin{problem}[\vfill\newpage]
      Let $R_0$ be the ratio of carbon-14 to carbon-12 in a living plant. Write a formula for the ratio of carbon-14 to carbon-12 in a plant $t$ years after it dies. (Remember that the half-life of this ratio is 5730 years.)
    \end{problem}

    \begin{solution} 
      $R_0 (\frac{1}{2})^{t / 5730}$
    \end{solution}

  \calcitem

    \begin{grading}
      2 points
    \end{grading}

    \begin{problem}[\vfill]
      When scientists used radiocarbon dating on the Winchester Table in 1976, they found that the ratio of carbon-14 to carbon-12 in the table was 91\% of the ratio of carbon-14 to carbon-12 in a sample of living wood. Based on this, which of the following choices describes the age of the table in 1976?
      \begin{enumerate}[nosep]
      \item Between 0 and 500 years.
      \item Between 500 and 1000 years.
      \item Between 1000 and 1500 years.
      \item More than 1500 years.      
      \end{enumerate}
      Explain your reasoning. 
    \end{problem}

    \begin{solution} 
      Using our formula we can fill in the following table
      \begin{center}
        \begin{tabular} {c| c} 
          $t$ & $R(t)$ \\ \hline
          0& $R_0$\\ 
          500& $.94R_0$\\ 
          1000  & $.88R_0$\\ 
          1500 &$.83 R_0$
        \end{tabular} \end{center} 
      The ratio of carbon-14 to carbon 12 would be $91 \%$ between 500 and 1000 years. 
    \end{solution}

  \item

    \begin{grading}
      0.5 points
    \end{grading}

    \begin{problem}[\nstretch{0.35}]
      Does it seem likely that the Winchester Table was King Arthur's Round Table from the 5th century? Why or why not?
    \end{problem}    

    \begin{solution} 
       No. 
    \end{solution}

  \end{enumerate} 

\item  

  \note{Preparation for thinking about the derivative of $e^{kx}$} 

  Here's the graph of an exponential function $f(x) = C \cdot b^x$. The tangent line at $x = 0$ is also shown, and it has slope $\frac{1}{3}$.
  \begin{center}
    \begin{tikzpicture}
      \begin{axis}[ytick={5, 10, 15, 20}, xtick={-20, -15, ..., 20}, gridgraph, ymin=0, ymax=20, axis equal image, width=4in]
        \addplot+[domain=-20:20]{5*e^(x/15)};
        \addplot+[domain=-8:8] {1/3*x+5}; 
      \end{axis}
    \end{tikzpicture}
  \end{center}

  \begin{enumerate}
  \item

    \begin{grading}
      1 point
    \end{grading}

    \begin{problem}[0.25in]
      What's $C$? How do you know? 
    \end{problem}

    \begin{solution}
      $C=5$. We know this because the $y$-intercept
      of the graph is $5$.
    \end{solution}

  \item

    \begin{grading}
      1 point
    \end{grading}

    \begin{problem}[\stretch{0.5}]
      Which of the following is true? How do you know?
      \[ b < 0 \hspace{1in} 0 < b < 1 \hspace{1in} b > 1\]
    \end{problem}

    \begin{solution}
      $b>1$. We know this because the function is increasing.
    \end{solution}

  \item
    \begin{grading}
      3.5 points. 
      \begin{itemize}[nosep]
      \item 0.5 for identifying this as a horizontal stretch
      \item 1 for having $y$-intercept 5
      \item 1 for looking reasonable when $x > 0$. (If $g$ isn't above $f$, they lose this point. If they've done some kind of horizontal stretch but it doesn't look close to a factor of $\frac{1}{2}$, deduct 0.5.)
      \item 1 for looking reasonable when $x < 0$. (Again, if $g$ isn't below $f$, they lose this point. If they've done some kind of horizontal stretch but it doesn't look close to a factor of $\frac{1}{2}$, deduct 0.5.) 
      \end{itemize}
    \end{grading}

    \begin{problem}[\stretch{0.75}]
      On the axes above, sketch the graph of $g(x) = f(2x)$. Explain in words how the graph of $g$ relates to the graph of $f$. 
    \end{problem}

    \begin{solution}
      The graph of $y=g(x)$ is the graph of $y=f(x)$ scaled horizontally by a factor
      of $\frac{1}{2}$.
      \begin{center}
        \begin{tikzpicture}
          \begin{axis}[ytick={5, 10, 15, 20}, xtick={-20, -15, ..., 20}, gridgraph, ymin=0, ymax=20, axis equal image, width=4in, restrict y to domain=0:20]
            \addplot+[domain=-20:20]{5*e^(x/15)};
            \addplot+[red, domain=-20:20]{5*e^(2*x/15)};
          \end{axis}
        \end{tikzpicture}
      \end{center}
    \end{solution}

  \item

    \begin{grading}
      0.5 points
    \end{grading}

    \begin{problem}
      On the graph above, sketch a line whose slope is $g'(0)$. 
    \end{problem}

    \begin{solution}
        \begin{center}
        \begin{tikzpicture}
          \begin{axis}[ytick={5, 10, 15, 20}, xtick={-20, -15, ..., 20}, gridgraph, ymin=0, ymax=20, axis equal image, width=4in, restrict y to domain=0:20]
            \addplot+[domain=-20:20]{5*e^(x/15)};
            \addplot+[red, domain=-20:20]{5*e^(2*x/15)};
            \addplot+[red, domain=-8:8] {2/3*x+5}; 
          \end{axis}
        \end{tikzpicture}
      \end{center}
    \end{solution}

  \item

    \begin{grading}
      1.5 points. This is exploratory, so they get full credit for any thoughtful attempt. 
    \end{grading}

    \begin{problem}[\vfill\newpage]
      (Exploratory) What is $g'(0)$? (You should be able to give a numerical answer.) Explain your reasoning. 
    \end{problem}

    \begin{solution}
      Since the graph of $y=g(x)$ is a horizontal scaling by a factor of $\frac{1}{2}$
      of the graph of $y=f(x)$, the tangent line to $y=g(x)$ at $x=0$ will also be
      a horizontal scaling by a factor of $\frac{1}{2}$ of the tangent line to $y=f(x)$
      at $x=0$. We are told that the slope of the tangent line to $y=f(x)$ at $x=0$
      is $\frac{1}{3}$, so $g'(0)$, the slope of the tangent line to $y=g(x)$ at $x=0$,
      is twice as steep, $\frac{2}{3}$.
    \end{solution}

  \end{enumerate}

\item  \begin{grading}
    2 points, full credit for any thoughtful answer
  \end{grading}

  \begin{problem}[\vfill]
    In workshop, you looked at $f(x) = b^x$ for some different values of $b$. In particular, you looked at approximations of $f'(x)$ for $x = 0, 1, 2, 3, 4, 5$. What patterns did you notice when you did this?     
  \end{problem}

  \begin{solution}
    See {{{{\#2}}(a)}} in workbook section 25!
  \end{solution}

  \probonly{We'll discuss these in class Friday!}

\end{enumerate}
\psreflection{For next class, read \S 9.4.}
\end{document}
